\chapter*{\textcolor{mycolor}{Introduction générale}}
\addcontentsline{toc}{chapter}{Introduction générale}

Un assistant personnel sait répondre à une question, programmer un rappel ou lire un message. Il ne fait rien de
tout cela si on ne le lui demande pas. Le problème est que les moments où l'on aurait besoin d'aide sont souvent
ceux où l'on n'y pense pas : le matin, au moment de partir, quand on referme la porte sans avoir pris son
portefeuille.

Le projet \textbf{LIFEWISE}, proposé par le Pr. Otman AGHZOUT dans le cadre du module de Computer Vision, part de ce
constat. Il s'agit de concevoir un assistant contextuel qui observe la situation réelle de la personne et qui
intervient de lui-même, au moment opportun. Pour cela, le projet associe trois domaines : l'Internet des objets,
pour capter l'environnement ; la vision par ordinateur, pour comprendre ce que voit une caméra ; et l'intelligence
artificielle embarquée (Edge AI), pour que tout se passe en local, sur un petit appareil, sans envoyer d'images à un
service distant.

Le sujet décrit un scénario global, à la maison, au travail et en déplacement. Nous ne le traitons pas en entier.
L'encadrant a fixé l'objectif de cette année : réussir de bout en bout un scénario pilote, l'assistance à la
sortie du domicile. Au moment où l'utilisateur part, le système vérifie que les objets qu'il emporte d'habitude,
le téléphone, les clés et le portefeuille, sont bien là. Si l'un d'eux n'est pas visible, le téléphone l'invite à
vérifier.

Cette formulation n'est pas un détail. Une caméra ne prouve qu'une chose : un objet n'est pas dans son champ. Elle
ne prouve pas qu'il a été oublié. Toute la chaîne est donc construite pour distinguer ce qui est observé de ce qui
est interprété, et l'alerte reste une invitation à vérifier.

La problématique peut se formuler ainsi : \textit{comment concevoir une chaîne locale, de la caméra jusqu'au
téléphone, qui détecte une situation précise et propose une aide utile au bon moment, sans affirmer plus que ce
que ses capteurs observent, et en respectant la vie privée de l'utilisateur ?}

Le projet est mené par treize élèves ingénieurs organisés en six sous-équipes, chacune responsable d'un maillon
de la chaîne. Ce rapport est structuré en cinq chapitres. Le premier présente le contexte général du projet,
l'analyse des besoins et l'organisation de l'équipe. Le deuxième dresse l'état de l'art des technologies
mobilisées. Le troisième décrit la conception et l'architecture du système. Le quatrième présente la réalisation
du prototype, module par module. Le cinquième rend compte des tests, de l'évaluation et de l'analyse des résultats.
Une conclusion générale et des perspectives terminent le rapport.
